\section*{Chapter \ref{ch:m1:us-equity-market-structure} --- US Equity Market Structure}

\begin{solution}{exo:m1:us-equity-market-structure:1}
D's bid (80 shares) and C's offer (50 shares) are \omterm{def:m1:us-equity-market-structure:lots}{odd lots} and are ignored.
Best bid 25.11 for $100 + 300 = 400$ shares (B, C); best offer 25.13 for 300
shares (B): $25.11 \times 25.13$, 400 by 300.
\end{solution}

\begin{solution}{exo:m1:us-equity-market-structure:2}
100 shares (\$4\,800); 40 (\$10\,040); 40 (\$39\,960); 10 (\$42\,000); 1
(\$712\,000).
\end{solution}

\begin{solution}{exo:m1:us-equity-market-structure:3}
Exchanges $49.4\% \times 17.6 = 8.7$ billion shares; principal dealers
$50.6\% \times 81.3\% \times 17.6 = 7.2$ billion; \omterm{def:m1:exchanges-brokers-venues:ats}{alternative trading systems}
1.7 billion. Average price $1.1\times10^{12}/17.6\times10^9 = \$62.50$.
\end{solution}

\begin{solution}{exo:m1:us-equity-market-structure:4}
B's protected offer of 25.13 for 300 shares was traded through. To buy 500 at
once: \omterm{def:m1:us-equity-market-structure:iso}{intermarket sweep orders} sent simultaneously for 300 at 25.13 to B and
for 200 at 25.14 to A. (C's 50 shares at 25.13 are unprotected; taking them
too is good execution, not an obligation.)
\end{solution}

\begin{solution}{exo:m1:us-equity-market-structure:5}
Maker-taker: taker pays 40.0030, \omterm{def:m1:us-equity-market-structure:makertaker}{maker} receives 40.0025, venue keeps 0.05
cent. Inverted: taker pays 39.9990, \omterm{def:m1:us-equity-market-structure:makertaker}{maker} receives 39.9986, venue keeps 0.04
cent. A \omterm{def:m1:us-equity-market-structure:makertaker}{maker} gives up $0.25 + 0.14 = 0.39$ cent a share to be hit first.
\end{solution}

\begin{solution}{exo:m1:us-equity-market-structure:6}
$\lambda\delta = 12 \times 600\times10^{-6} = 0.72\%$, that is 168 seconds a
day. When changes cluster, several fall within one interval of length
$\delta$ and the stale intervals overlap: their union is shorter than their
sum.
\end{solution}

\begin{solution}{exo:m1:us-equity-market-structure:7}
0.255\% in both cases. The tape's staleness relative to a \omterm{def:m1:us-equity-market-structure:sip}{direct feed} depends
only on the extra delay $\delta$: both views wait for the venues to reprice.
The venues' delays control something else --- how long the \emph{quotes
themselves} lag the fair price, which is the window in which they can be
sniped.
\end{solution}

\begin{solution}{exo:m1:us-equity-market-structure:8}
(i) The benchmark is the tape's NBBO, stale exactly when prices are moving:
an execution ``at the NBBO'' during those intervals is worse than the true
best price. (ii) Timestamping on receipt ignores the time the order then
spent at the wholesaler or in routing; the fair comparison is with the NBBO
at execution, and clock differences of a millisecond change the answer.
(iii) In high-priced stocks the best prices were \omterm{def:m1:us-equity-market-structure:lots}{odd lots}, excluded from the
NBBO: the benchmark was wider than the real market, so ``at or better'' was
easy. More generally, ``at or better than'' says nothing about size or about
the midpoint, which is what the order could often have obtained.
\end{solution}

\begin{solution}{pb:m1:us-equity-market-structure:1}
\textbf{1.} $\lambda\delta = (8/60) \times 900\times10^{-6} = 0.012\%$.
\textbf{2.} 2.8 seconds.
\textbf{3.} The fastest venue's, after 60 microseconds: a higher bid
anywhere raises the best bid.
\textbf{4.} Only when the \emph{slowest} venue reprices, after 500
microseconds: the best offer is the lowest, and the stale low offers remain
the best until the last of them is gone.
\textbf{5.} All four: the fastest reprices at 60 microseconds and the trader
arrives at 40.
\textbf{6.} Old offer $m+1$, new fair value $m+J$: profit $J-1$ ticks, that
is 0, 1 and 2 cents.
\textbf{7.} $0.25 \times 1 + 0.10 \times 2 = 0.45$ cent.
\textbf{8.} $0.45 \times 4 \times 300/100 = \$5.40$ per jump; with 3\,120
jumps a day, \$16\,848.
\textbf{9.} $0.25\times0.70 + 0.10\times1.70 = 0.345$ cent: \$12\,917 a day.
One-tick jumps lose the fee and are left alone; only jumps of two ticks or
more are taken.
\textbf{10.} $0.45 \times 3 \times 3\,120 = \$4\,212$ a day, 1.05 cent per
share traded.
\textbf{11.} No: $0.35 - 1.05 = -0.70$ cent a share. It earns \$1\,400 a day
and loses \$4\,212.
\textbf{12.} It would save \$1.06 million a year for \$2.5 million: not on
these numbers. Its real choices are to quote wider, to quote smaller, or to
stop.
\textbf{13.} Each would then again be slower than whoever invests next; the
trader upgrades too, the sniping profit is unchanged or merely reallocated,
and the industry has spent \$10 million a year for nothing: an arms race, in
which speed is valuable only relative to others.
\textbf{14.} A quote is safe if its \omterm{def:m1:what-a-trading-firm-does:market-maker}{market maker} reprices before a delayed
order can arrive, that is within $40 + 350 = 390$ microseconds: the three
fastest. The slowest, at 500 microseconds, is still exposed.
\textbf{15.} To a venue still showing the old, lower offer; the order either
arrives after the quote has been repriced or taken and is not filled, or
executes against a fresh higher offer. The client chases the price.
\textbf{16.} No. Rule 611 compares the execution with protected quotations
as displayed at the time; every venue's displayed offer was at or above the
price paid. The rule protects displayed prices, not fair ones, and says
nothing about stale quotes.
\textbf{17.} $0.25\times0.90 + 0.10\times1.90 = 0.415$ cent: a lower fee cap
makes sniping \emph{more} profitable.
\textbf{18.} No: the opportunity comes from quotes lagging the fair price,
not from routing obligations. Rescission would change who must route where,
not the race to a stale quote.
\textbf{19.} \textbf{0.345 cent per share.}
\textbf{20.} The investors who trade with the \omterm{def:m1:what-a-trading-firm-does:market-maker}{market makers}, through the
wider spreads and smaller sizes that pay for the losses of question 10.
\end{solution}

\begin{solution}{iq:m1:us-equity-market-structure:1}
The highest protected bid and lowest protected offer across exchanges, with
sizes. Officially the SIPs compute it; every serious firm computes its own
from \omterm{def:m1:us-equity-market-structure:sip}{direct feeds}. They disagree because each observer receives each venue's
update after a different delay and with a different clock: the NBBO is
defined relative to a place and a time, and ``the NBBO at 10:31:04.123456''
has no unique answer.
\iqlookfor{relativity of the consolidated view, not just ``SIP is slow''.}
\end{solution}

\begin{solution}{iq:m1:us-equity-market-structure:2}
A limit order flagged to tell the receiving venue to execute without
checking away markets, because the sender has simultaneously sent orders to
clear every better-priced protected quote. It exists so that an order larger
than the best displayed size can take several price levels on several venues
in one instant instead of being rerouted level by level while the market
moves away.
\iqlookfor{responsibility for \omterm{def:m1:us-equity-market-structure:opr}{trade-through} compliance shifting to the
sender.}
\end{solution}

\begin{solution}{iq:m1:us-equity-market-structure:3}
Takers hit the \omterm{def:m1:us-equity-market-structure:makertaker}{inverted venue} first because it pays them, so a resting order
there executes before the same price on any maker-taker venue: the fee buys
time priority across venues. It is worth paying when the queue at the best
price is long and the spread is pinned at one tick, so that priority, not
price, decides who trades --- and the fills obtained early, before the queue
has been eaten, are also the less adversely selected ones.
\iqlookfor{queue position as the thing being bought.}
\end{solution}

\begin{solution}{iq:m1:us-equity-market-structure:4}
Per symbol, a small fixed array of per-venue best quotes (prices as
integers) plus the cached best bid and offer and venue bit masks; an update
touches one slot and, only if it was at the best or improves it, a rescan of
sixteen entries that fits in a cache line or two. Symbols in a dense array
indexed by an integer identifier, not a hash of a string. One thread per feed
decodes and writes; shard symbols across book-building threads so that each
symbol has a single writer and no locks; publish to strategies through a
sequence-locked snapshot or a ring buffer. Pin threads, pre-allocate
everything, keep the hot structure contiguous.
\iqlookfor{single-writer sharding and cache-aware layout.}
\end{solution}

\begin{solution}{iq:m1:us-equity-market-structure:5}
For retail investors mostly not: they get small price improvement and no
commission. For institutions, \omterm{def:m1:us-equity-market-structure:dark}{dark pools} reduce information leakage. The
concern is for price formation: off-exchange trades use the NBBO without
contributing to it, and if the least-informed flow is filtered away from
exchanges, displayed spreads widen for everyone who must trade there, and
the benchmark for ``improvement'' degrades. It is also a competition question:
a few wholesalers see, and price, most retail flow.
\iqlookfor{segmentation and \omterm{def:m1:what-a-trading-firm-does:adverse-selection}{adverse selection}, with both sides.}
\end{solution}

\begin{solution}{iq:m1:us-equity-market-structure:6}
(i) Venues need no longer route away or reject orders that would trade
through: small exchanges that survive on protected-quote status lose flow,
and volume concentrates; position connectivity and quoting capital on the
venues likely to gain. (ii) Best execution becomes the broker's judgement
rather than a mechanical rule: brokers' routers and their transaction-cost
analysis matter more, so a \omterm{def:m1:what-a-trading-firm-does:market-maker}{market maker} should expect more scrutiny of fill
quality by venue and compete on it. (iii) Locked and crossed markets become
legal: displayed prices across venues will overlap briefly, creating
arbitrage for the fastest and new states that every quoting engine and risk
check must handle. Also likely: more differentiated venue designs (delays,
batch auctions) no longer constrained by protected-quote requirements.
\iqlookfor{second-order effects and concrete firm actions, not a list of
opinions.}
\end{solution}
